\subsubsection{终端命题：黄金分支时间、6D 微态与 Zeckendorf--GUT 的离散三分支}
\label{subsubsec:window6-terminal-propositions}

依 \cite{Internal2026POMGenerativeAddressThreeRegisterBridgeRH}（2026-02-09 版本）所固定的协议口径，本文将“时间”视为\emph{单生成扫描算子}的迭代计数，并将黄金分支 \(\alpha=\varphi^{-1}\) 视为在 Hurwitz/Markov 反共振基线下的规范代表（\S\ref{sec:principles}）。
在此框架内，window-\(6\) 作为最小锚点同时控制：二元窗读出的离散时间刻度、6D 超立方体微态到稳定语法的压缩、以及由 Zeckendorf--Fibonacci 边界塔诱导的统一维数分支结构。
本小节仅陈述终端断言与可审计接口指针，不叙述中间推导。

\begin{enumerate}
  \item \textbf{时间的本体刻画并非“旋转本身”，而是“单生成扫描”的迭代参数；黄金旋转仅为其最小接口表示。}

  \begin{theorem}[扫描迭代时间与黄金分支的规范性]
  \label{thm:terminal-time-protocol-golden-branch}
  在扫描—投影读出框架中，给定单生成作用 \(T^t\) 与二元窗读出
  \(s_t(x)=\ind_W(T^t x)\in\{0,1\}\)，时间参数 \(t\) 被定义为 \(T\) 的迭代计数，而非外加连续坐标（\S\ref{sec:spg}；并见 \S\ref{subsec:time-protocol-entropy-vs-rotation}）。
  在对象层，若存在最大等距因子，则可将扫描相位表示为紧交换群上的 Kronecker 平移；一维圆周情形下即为无理旋转与区间窗的符号化接口（附录 \ref{app:cut_project_background} 的 Sturmian 原型与 phason 参数化注 \ref{rem:app-cutproject-phason-time}）。
  在协议层，本文采用 Zeckendorf/Ostrowski 兼容的规范归一化并以 Hurwitz/Markov 的最坏逼近尺度固定反共振基线，从而将黄金分支 \(\alpha=\varphi^{-1}\)（及其 \(\mathrm{SL}_2(\ZZ)\) 等价类）区分为该口径下的规范代表（定义 \ref{def:markov-constant} 与 \S\ref{sec:principles}）。
  \end{theorem}

  \begin{corollary}[\texorpdfstring{$\log\varphi$}{log phi} 刻度锁定与内生相关时间闭式]
  \label{cor:terminal-tau-corr-logphi}
  在稳定语法取黄金均值约束并以 Parry 最大熵平衡态作为协议基准时，时间单位由熵率 \(\log\varphi\) 锁定（命题 \ref{prop:time-unit-logphi-protocol}）。
  同时，在真实输入 \(40\) 状态核的 \(\chi\) 切片 \(t=0\) 的 Parry 平衡态下，谱隙比满足
  \(
  \rho_{\mathrm{corr}}(0)=\varphi^{-1}
  \)
  且内生相关时间具有闭式
  \[
  \tau_{\mathrm{corr}}(0)=\frac{1}{-\log\rho_{\mathrm{corr}}(0)}=\frac{1}{\log\varphi}
  \]
  （注 \ref{cor:chi-slice-t0-golden-tau}）。
  该框架亦给出将步数时间重标定为事件时间的可审计换元（\S\ref{subsec:chi-ldp} 中关于事件时钟的表格证书）。
  \end{corollary}

  \item \textbf{6D 超立方体微态与 window-\(6\) 折叠：从 \(\Omega_6=\{0,1\}^6\) 到稳定类型 \(X_6\) 的不可约压缩。}

  \begin{theorem}[window-\(6\) 的 dyadic 折叠：\texorpdfstring{$64\to 21$}{64->21} 与刚性退化直方图]
  \label{thm:terminal-foldbin6-64-to-21-hist}
  在最小分辨率 \(m=6\) 下，窗口微态空间为 \(\Omega_6=\{0,1\}^6\)（\(|\Omega_6|=64\)），稳定语法空间为 \(X_6\subseteq\Omega_6\)（禁止相邻 \(11\)，故 \(|X_6|=F_8=21\)，见 \S\ref{sec:folding}）。
  二进制区间折叠 \(\mathrm{Fold}^{\mathrm{bin}}_6:\{0,\dots,63\}\twoheadrightarrow X_6\)（定义 \ref{def:fold-bin}）的纤维大小仅取
  \(\{2,3,4\}\)，并且退化度直方图为
  \[
  \#\{d^{\mathrm{bin}}_6=2\}=8,\qquad
  \#\{d^{\mathrm{bin}}_6=3\}=4,\qquad
  \#\{d^{\mathrm{bin}}_6=4\}=9
  \]
  （推论 \ref{cor:fold-bin-m6-fiber-hist}；并见表 \ref{tab:fold_bin_gauge_m6_invariants}）。
  \end{theorem}

  \begin{proposition}[6D 可达均值凸体与可微分指纹（压力梯度/Hessian）]
  \label{prop:terminal-6d-reachable-mean-convex-body}
  在 \S\ref{subsec:pom-potential-convex-geometry} 的势向量口径下，选取 \(6\) 个二元通道作为势分量时，长度 \(n\) 的轨道产生计数向量
  \(N_n\in\{0,\dots,n\}^6\)，归一化后落在 \([0,1]^6\)。
  相应的可达均值集合 \(\mathcal{R}\subset[0,1]^6\) 为紧凸体；其 \(3\) 维切片（固定三个坐标或三个线性组合）构成“6D 立方体的 3D 切面”。
  更重要的是，压力函数 \(P(\theta)=\log\rho(M_\theta)\) 的梯度与 Hessian 在切片上的读出给出该凸体的局部一阶/二阶指纹，并把几何以凸对偶形式内生化为可审计对象（\S\ref{subsec:pom-potential-convex-geometry}）。
  \end{proposition}

  \item \textbf{window-\(6\) uplift 的三项尾部偏移 \(\{21,34,55\}\) 与 Zeckendorf--GUT 顶层项的刚性对齐。}

  \begin{theorem}[window-\(6\) uplift 的离散三分支与 Zeckendorf--GUT 顶层对齐]
  \label{thm:terminal-window6-tail-three-branch}
  对任意 \(w\in X_6\)，其 dyadic uplift 预像相对于 Zeckendorf 值 \(V_6(w)\) 的非零偏移仅可能取自
  \[
  \Delta_{\mathrm{tail}}(6)\subseteq\{F_8,F_9,F_{10}\}=\{21,34,55\}
  \]
  （推论 \ref{cor:fold-bin-m6-fiber-hist}；推论 \ref{cor:fold6-tail-offsets-gut-top-terms}）。
  另一方面，边界词塔满足 \(|X_m^{\mathrm{bdry}}|=F_{m-2}\)，从而对任意候选李代数 \(\mathfrak g\)（\(D=\dim\mathfrak g\)）的 Zeckendorf 分解诱导唯一的分辨率签名 \(\Sigma(\mathfrak g)\)，并给出
  \(D=\sum_{m\in\Sigma(\mathfrak g)}|X_m^{\mathrm{bdry}}|\)（\S\ref{subsec:bdry-tower-zeck-gut}；表 \ref{tab:zeckendorf_gut_signatures}）。
  在该字典下，\((SU(5),SO(10),E_6)\) 的 Zeckendorf--GUT 分解最高 Fibonacci 项分别为 \((F_8,F_9,F_{10})\)，从而与 window-\(6\) uplift 的三项尾部偏移严格对齐，并在协议层导出一条离散三分支选择通道。
  \end{theorem}

  \item \textbf{家族数的分支锁定：\texorpdfstring{$N_f=3$}{Nf=3} 在 window-\(6\) 精确选中 \texorpdfstring{$F_9=34$}{F9=34} 分支。}

  \begin{theorem}[家族数锁定的三分支选择律]
  \label{thm:terminal-family-uplift-lock}
  以每代 \(SU(5)\) 型手征物质标准计数 \(15\) 为外部输入，令总计数 \(M(N_f)=15N_f\)。
  则对 \(N_f\in\{2,3,4\}\)，\(M(N_f)\) 的 Zeckendorf 分解最高项依次为 \(F_8,F_9,F_{10}\)；
  因而家族数在 window-\(6\) 层级被提升为离散分支选择量，且 \(N_f=3\) 精确选中 \(F_9=34\) 分支并与 \(SO(10)\) 顶层项对齐（表 \ref{tab:window6_family_uplift_lock}）。
  \end{theorem}

  \item \textbf{window-\(6\) 内部的 \texorpdfstring{$1\oplus 8\oplus 12$}{1+8+12} 离散劈裂与 \texorpdfstring{$\delta=34$}{δ=34} sheet parity 选择律。}

  \begin{proposition}[\texorpdfstring{$1\oplus 8\oplus 12$}{1+8+12} 劈裂与轻/重字典的可枚举性]
  \label{prop:terminal-window6-1-8-12-split}
  在 dyadic 折叠口径下，uplift-choice 阈值导致 \(\Xi(w):=|(\mathrm{Fold}^{\mathrm{bin}}_6)^{-1}(w)|\in\{2,3,4\}\) 的三值塌缩，并据此在 \(X_6\) 上诱导一个完全可枚举的三扇区分解
  \[
  X_6=X_6^{\mathrm{zero}}\ \oplus\ X_6^{\mathrm{light}}\ \oplus\ X_6^{\mathrm{heavy}},
  \qquad
  |X_6^{\mathrm{zero}}|=1,\ |X_6^{\mathrm{light}}|=8,\ |X_6^{\mathrm{heavy}}|=12,
  \]
  其完整词表与阈值条件由表 \ref{tab:fold6_bin_uplift_choice_collapse} 给出。
  通过四步平移 uplift 同构 \(X_n\simeq X_{n+4}^{\mathrm{bdry}}\)（定理 \ref{thm:boundary-shift4-uplift-isomorphism}），该离散字典可规范抬回 \(m=10\) 的边界层，从而将 “\(21-(8+1)=12\)” 从计数目标提升为可逐项核验的轻/重字典接口（见 \S\ref{subsec:bdry-tower-zeck-gut} 的相应段落）。
  \end{proposition}

  \begin{corollary}[\texorpdfstring{$\delta=34$}{δ=34} 的二层覆盖与离散奇偶选择定则]
  \label{cor:terminal-delta-parity-rule}
  window-\(6\) 的边界扇区 \(X_6^{\mathrm{bdry}}\) 在 dyadic uplift 下呈现规范的二层覆盖：每个边界词恰有两条预像，且两者差值严格为 \(34=F_9\)，从而诱导一条内生的 \(\ZZ_2\) sheet parity 分级（表 \ref{tab:fold6_boundary_sheet_pairs}）。
  在要求可观测耦合与该分级相容的协议门口径下，bulk 与 boundary 的混合项必须经由 \(\delta=34\) 的 sheet 翻转实现，从而给出无需连续微调的离散选择律（参见 \S\ref{subsec:bdry-tower-zeck-gut} 中的“\(\delta\)-奇偶”讨论段落）。
  \end{corollary}

  \item \textbf{多重度极值与饱和判别：dyadic 分辨率 \texorpdfstring{$2^m$}{2^m} 与 Fibonacci 分层的指数级丢失结构。}

  \begin{theorem}[最大纤维、组合上界与饱和的 Diophantine 刚性]
  \label{thm:terminal-foldbin-saturation-diophantine}
  对任意 \(m\ge 1\)，二进制区间折叠 \(\mathrm{Fold}^{\mathrm{bin}}_m\) 的最大纤维由 \(0^m\) 取得，并满足显式组合上界与饱和判别式（推论 \ref{cor:fold-bin-saturation}）。
  饱和条件强迫 \(2^m\) 与某个 Fibonacci 数发生指数级逼近，并可改写为对数线性形式的指数小上界；结合 Baker--W\"ustholz/Matveev 型下界得到“仅有限次饱和”的机制解释（注 \ref{rem:fold-bin-finite}）。
  \end{theorem}

  \item \textbf{切投影与折叠的统一接口：稳定语法直方图作为跨分辨率不变量。}

  \begin{proposition}[cut-and-project \texorpdfstring{$\Rightarrow$}{=>} 折叠直方图不变量的接口]
  \label{prop:terminal-cut-project-to-fold-hist}
  切投影/模型集的局部贴片可通过内空间窗集 \(W\) 的指示函数作滑动窗口编码；在此基础上施加黄金均值约束并施行 \(\Fold_m\) 折叠，即可把高维几何数据在固定分辨率下压缩为稳定语法类型直方图。
  该直方图的递归尺度由 \(\varphi\) 控制，因而适合作为跨分辨率比较的无量纲不变量（附录 \ref{sec:high-dimensional-cut-project} 与附录 \ref{app:cut_project_background}）。
  \end{proposition}

  \begin{corollary}[环面参数化与旋转扫描的可审计原型]
  \label{cor:terminal-rotation-window-prototype}
  在一维原型中，环面旋转 \(T(x)=x+\alpha\pmod 1\) 与区间窗 \(W=[0,\beta)\) 的读出 \(s_t(x)=\ind_W(T^t x)\) 给出“旋转+\(\)窗”的标准符号接口；有限窗口 \(m\) 的局部块统计对应于模型集的有限贴片计数，从而为“扫描—窗口读出—折叠稳定语法”的统一几何语境提供可审计基准（附录 \ref{app:cut_project_background}）。
  \end{corollary}

  \item \textbf{与 6D 超立方体直接相关的统一性收束：黄金时间与 Zeckendorf 边界塔共同导出离散统一分支。}
  \end{enumerate}

\begin{theorem}[6D 微态—黄金时间—GUT 分支的离散统一]
\label{thm:terminal-6d-microstate-golden-time-gut-branch}
在最小局部分辨率 \(m=6\) 下，6D 超立方体微态 \(\Omega_6\) 在 dyadic 折叠口径中经 \(\mathrm{Fold}^{\mathrm{bin}}_6\) 压缩到稳定语法 \(X_6\)，并在 uplift 层面仅允许三种非零尾部偏移 \(\{F_8,F_9,F_{10}\}=\{21,34,55\}\)，从而在 window-\(6\) 层级导出离散三分支（定理 \ref{thm:terminal-window6-tail-three-branch}）。
与此同时，Zeckendorf--Fibonacci 边界塔签名将李代数维数 \(D=\dim\mathfrak g\) 唯一分解为边界层计数之和，并使 \((SU(5),SO(10),E_6)\) 的顶层项与 \((F_8,F_9,F_{10})\) 精确对齐（\S\ref{subsec:bdry-tower-zeck-gut}；表 \ref{tab:zeckendorf_gut_signatures}）。
进一步地，家族数 \(N_f\) 通过 \(15N_f\) 的顶层 Zeckendorf 项在 \(N_f\in\{2,3,4\}\) 上实现分支锁定，尤其 \(N_f=3\) 离散选中 \(F_9=34\) 分支并与 \(SO(10)\) 顶层对齐（定理 \ref{thm:terminal-family-uplift-lock}）。
因此，在该协议的最小分辨率锚点上，“统一候选的顶层选择”呈现为一个由\emph{内生时间刻度（\(\log\varphi\)）}与\emph{dyadic uplift 尾部偏移}共同决定的离散选择问题，而非连续可调参数族。
\end{theorem}

\paragraph{材料指针（可审计入口）}
\begin{itemize}
  \item \textbf{时间=迭代计数（对象层/协议层拆分）与 \(\log\varphi\) 刻度：}
  \S\ref{sec:spg}；\S\ref{subsec:time-protocol-entropy-vs-rotation}；命题 \ref{prop:time-unit-logphi-protocol}。
  \item \textbf{黄金点相关时间闭式与事件换时：}
  注 \ref{cor:chi-slice-t0-golden-tau}；以及 \S\ref{subsec:chi-ldp} 中的事件时钟表格（例如 \texttt{tab\_real\_input\_40\_event\_time\_dilation} 等生成片段）。
  \item \textbf{window-\(6\) dyadic 折叠的 \((2,3,4)\) 退化谱直方图与 gauge 不变量：}
  推论 \ref{cor:fold-bin-m6-fiber-hist}；表 \ref{tab:fold_bin_gauge_m6_invariants}；脚本 \texttt{scripts/exp\_fold\_bin\_gauge\_m6\_invariants.py}。
  \item \textbf{Zeckendorf--GUT 签名与 window-\(6\) 三分支对齐：}
  \S\ref{subsec:bdry-tower-zeck-gut}；表 \ref{tab:zeckendorf_gut_signatures}；脚本 \texttt{scripts/exp\_boundary\_tower\_gut\_signatures.py}。
  \item \textbf{家族数锁定（\(15N_f\) 顶层项）与 window-\(6\) uplift 对齐：}
  表 \ref{tab:window6_family_uplift_lock}；脚本 \texttt{scripts/exp\_window6\_family\_uplift\_lock.py}。
  \item \textbf{\(\delta=34\) 二层覆盖与 sheet parity：}
  表 \ref{tab:fold6_boundary_sheet_pairs}；表 \ref{tab:fold6_bin_uplift_choice_collapse}。
  \item \textbf{最大纤维饱和判别与 Diophantine 刚性：}
  推论 \ref{cor:fold-bin-saturation} 与注 \ref{rem:fold-bin-finite}（附录 \ref{app:fold-multiplicity}）。
  \item \textbf{cut-and-project 与“旋转+\(\)窗”原型、phason 时间定位：}
  附录 \ref{sec:high-dimensional-cut-project}；附录 \ref{app:cut_project_background}（注 \ref{rem:app-cutproject-phason-time}）。
\end{itemize}

